% The node as modules. Three rows, explicit coordinates, and every dependency
% drawn as a direct edge: no right-angle routing, so no line can pass through a
% box that is not its target.
\begin{tikzpicture}[font=\sffamily\scriptsize, x=1cm, y=1cm]

\newcommand{\umlbox}[4]{%
  \node[umlclass, text width=32mm] (#1) at (#2) {%
    \begin{minipage}{32mm}
      \centering\bfseries #3\par\vspace{1pt}
      \hrule\vspace{2pt}
      \raggedright\normalfont #4\vspace{1pt}
    \end{minipage}};}

% ---------------- row 1: the three modules this chapter writes
\umlbox{ident}{-4.6,3.4}{identity}{name, node\_id, version\\ built, uid[3]\\[1pt]
  print() at boot}
\umlbox{ind}{0,3.4}{indicator}{set(state)\\[1pt]
  {\itshape six states, two of them\\ unreachable until chapter 18}}
\umlbox{selfck}{4.6,3.4}{selfcheck}{clock, pll source\\ flash latency\\[1pt]
  {\itshape fails loudly, never warns}}

% ---------------- row 2: the one state, and what reads and writes it
\umlbox{ctrl}{-4.6,0.7}{control}{period, loop, plant\\[1pt]
  {\itshape chapters 2, 3 and 7}}
\umlbox{state}{0,0.7}{node\_state\_t}{position, velocity\\ effort\_cmd, effort\_est\\
  mode, faults, seq\\[1pt] {\itshape one instance, one owner}}
\umlbox{bus}{4.6,0.7}{bus}{controller, protocol\\ membership\\[1pt]
  {\itshape chapters 9 to 13}}

% ---------------- row 3: what sits above, and what sits below
\umlbox{mw}{-4.6,-2.1}{middleware}{client, joint messages\\[1pt] {\itshape chapters 14 and 15}}
\umlbox{bsp}{0,-2.1}{bsp}{startup, system, uart\\ board.h\\[1pt]
  {\itshape chapter 4 turns this into\\ a board file to hand over}}
\umlbox{safety}{4.6,-2.1}{safety}{safe states, latching\\[1pt] {\itshape chapter 18}}

% ---------------- dependencies, every one a direct edge
\draw[flow] (ident)  -- (state);
\draw[flow] (ind)    -- (state);
\draw[flow] (selfck) -- (state);
\draw[flow] (ctrl)   -- node[lbl]{writes} (state);
\draw[flow] (bus)    -- node[lbl]{reads} (state);
\draw[flow] (mw)     -- (ctrl);
\draw[flow] (state)  -- (bsp);
\draw[flow] (ctrl)   -- (bsp);
\draw[flow] (bus)    -- (bsp);
\draw[faultedge] (safety) -- node[lbl]{may force} (bus);

\node[umlnote, text width=52mm, anchor=north west] at (-6.3,-3.6)
  {Nothing points upward. The control loop does not know a bus exists: it
   produces a state and consumes a command, and something above it decides where
   those go. That is what lets chapter 16 close a loop over the network without
   touching chapter 2, and what lets the whole node be tested on a host machine
   in chapter 20.};
\node[umlnote, text width=40mm, anchor=north east] at (6.3,-3.6)
  {One state structure, one owner. The most common way a machine like this
   becomes untrustworthy is two modules each holding a slightly different idea
   of where the axis is.};
\end{tikzpicture}
